\minitopic{形式化首先是一项表示选择}把信念写成概率并不只是把模糊语言换成数字。研究者必须先选择命题空间：疾病是否只分“有／无”，还是包含不同亚型；天气是否只分“下雨／不下雨”，还是还要区分时间和区域。过粗空间会隐藏行动相关差异，过细空间又需要无法支持的参数。随后还要选择参照类、时间窗口和条件变量。任何计算在这些选择之后才开始，因此形式结果的严格性不能反向证明表示本身正确。

\minitopic{概率公理约束一致性而不提供全部内容}若主体对互斥完备命题给出的置信度总和不为一，赌局或评分规则可以揭示某种内部不一致。恰当评分规则进一步说明，诚实报告真实置信度具有期望优势\parencite{joyce1998}。然而公理没有告诉我们初始置信度应是多少，也没有判断模型遗漏了哪项假设。一个内部一致的阴谋论者可以给所有事件分配概率；问题在于他的似然、来源和更新是否与世界适当连接。形式规范因此是必要约束，却不是认识充分条件。

\minitopic{基础率是关于参照类的主张}开篇检测只有在患病率真正适用于当前患者时才可直接代入。年龄、症状、暴露史和就诊选择可能使患者属于不同参照类。忽略基础率会夸大阳性证据，机械使用一般人群基础率也可能低估个案信息。研究者应说明基础率来自何时、何地、哪类人群，样本如何进入数据；若多个参照类给出不同先验，应做敏感性分析，而不是把其中一个数字包装成客观起点。

\minitopic{似然衡量区分力而非结论概率}证据$E$在假设$H$下常见，只说明$P(E\mid H)$高；若$E$在非$H$下同样常见，它几乎没有区分力。贝叶斯因子比较不同假设对同一证据的预测，避免把“假设能解释证据”误作“证据强烈支持假设”。但是似然往往来自模型估计而非直接观察，模型误设、测量误差和选择性报告都会扭曲比值。对极大似然比尤其应追问：分母是否被低估，竞争假设是否遗漏，证据是否在看过结果后才被选择。

\begin{argumentbox}
报告一次概率更新时，不只写后验数值；同时写明命题空间、基础率来源、似然来源、独立性假设、遗漏假设和敏感性范围。若改变一个合理设定便使结论反转，最重要结果不是某个点估计，而是模型尚不稳健。
\end{argumentbox}

\minitopic{证据重复计算常藏在数据谱系中}两家网站引用同一报告、两次检测使用同一污染样本、三个模型在同一训练数据上给出相似结果，都可能产生条件相关。给定目标假设后，若一项结果出现会提高另一项出现的概率，就不能简单相乘似然。相关结构可以用联合概率、层级模型或显式因果图表示；即使无法精确估计，也应做上下界或情景分析。把来源谱系带进形式模型，是避免“看似证据很多、实际只有一条链”的关键。

\minitopic{校准与分辨力回答不同问题}总是报告基础率的系统可能完美校准，却无法区分个案；极具分辨力的系统也可能系统性过度自信。校准要按足够大样本和适当时间窗口检验，分辨力则比较发生与未发生事件的概率分布。分组校准还可能与其他公平目标冲突。因而“模型准确”至少应拆成总体错误、概率校准、排序能力、分组表现和极端情形稳健性。单一平均指标不能承担全部认识评价。

\minitopic{不确定性也有不同来源}随机不确定性来自事件本身的变异，认识不确定性来自数据不足、参数不明或模型竞争；分布变化则使过去估计不再适用于当前环境。三者需要不同回应：增加样本未必减少不可消除的随机性，精确计算也不能修复模型类别错误，重新校准则可能应对环境变化。把所有不确定性压成一个置信区间会隐藏应采取的行动。形式认识论的任务之一，正是把“不知道”细分为可通过什么证据减少、无法通过当前方法减少以及尚未被模型表达的部分。

\minitopic{置信度的解释会改变规范问题}把数值理解为长期频率、主体的主观信念程度，还是在公平赌局中愿意接受的比率，会导向不同检验。一次事件没有可直接观察的长期频率，纯主观解释又可能使任意先验显得宽容；行为解释则会受风险厌恶和财富变化干扰。实际研究常把这些解释组合使用：概率由群体频率与模型产生，再作为个体置信输入，最后通过选择行为间接表现。组合并非错误，但必须说明每次转换需要哪些桥梁，否则一个统计总体中的比例会被悄悄当成特定主体应有的信心。

\minitopic{先验不是更新前可以遗忘的背景}两个研究者使用相同证据而得到不同后验，分歧可能来自先验。先验可以编码既有研究、机制知识与基础率，也可能携带传统偏见、方便假设或任意平滑。称其“主观”不能终止评价：应追问它对哪些人群和时间适用，是否把争议结论预先写入模型，以及数据量增加后后验多快摆脱它。报告多组可辩护先验及其结果，比假装存在一个无立场起点更能显示结论强度。若结论只在极窄先验下成立，真正需要辩护的就是先验而非计算。

\minitopic{开放世界会使完备假设空间失效}贝叶斯归一化通常假定列出的假设互斥且完备，但现实调查经常遗漏未知故障、混合原因或尚未发现的疾病亚型。若模型只比较“机器正常”和“传感器损坏”，来自软件更新的第三种错误会被错误分配给前两项，后验仍然精确却回答了封闭世界中的问题。研究者可加入“其他原因”类别、使用异常检测、追踪残差，或明确将结论限定为“在已列模型中”。保留未知项会降低表面确定性，却防止模型以错误的完备性垄断解释。

\minitopic{证据选择本身需要进入模型}若研究者在看到结果后才决定报告哪项指标、何时停止实验或把哪个亚群当作目标，普通条件化无法自动纠正选择偏差。观察到“有一项检验显著”与事先指定的那项检验显著，似然不同；连续查看数据直到跨过阈值，也会增加偶然成功。完整证据不只包括最终数字，还包括采样规则、停止规则、排除标准和未报告比较。形式化的优势正是可以把这些程序写入概率模型，而不是把模型用于最终结果却把结果如何被挑选留在模型外。

\minitopic{证据关系不等于因果关系}变量之间的条件相关可以提高预测，却不说明干预其中一个会改变另一个。医院数据中，接受高强度治疗者死亡率较高，可能因为病情严重者更常接受治疗；把相关性直接作为治疗有害的证据，会混淆选择机制。认识论上，“支持某预测”与“支持某解释或干预”是不同任务。应先说明问题是预测、诊断还是因果决策，再选择相应模型与数据。一个在预测上校准良好的系统，仍可能在回答“若采取行动会怎样”时系统错误。

\minitopic{区间与集合有时比单一概率更诚实}当样本有限、先验争议大或多个模型难以排序时，给出一个精确到小数点后的概率会制造不存在的信息。可以用概率区间、模型集合或多种情景表达“不确定到什么程度”。代价是部分命题无法得到唯一排序，某些决定也不会由计算自动确定；收益则是把证据不足与立场分歧保留下来。区间不应无限宽到失去约束，而应由可说明的样本界限、模型差异和参数范围产生，并随新证据收窄或移动。

\begin{argumentbox}
一份可审计的形式分析至少回答六问：研究的命题空间是否适合任务；先验与基础率从何而来；证据选择过程怎样形成；哪些依赖与遗漏被模型表达；合理替代模型是否改变结论；数值精度是否超过材料能够支持的程度。公式正确只能通过其中一问。
\end{argumentbox}
